\section{Kinematics and Gait Planning}

\subsection{Forward and Inverse Kinematics}

To accurately mandate the position of a limb's endpoint
$(x, y, z)$ relative to the geometric origin of the base frame,
a non-linear inverse kinematics (IK) model derives the
necessary articular angles $( \theta_1, \theta_2, \theta_3 )$
for the hip, thigh, and calf joints.
The physical topology structuring the Cartesian workspace includes
the platform length $L$, platform width $W$, and individual
linkage spans $l_1, l_2, $ and $l_3$.

The generalized mathematical solution avoids geometric singularities
by enforcing strictly restricted bounding limits calculated
dynamically during system boot initialization.

\begin{figure}[htbp]
    \centering
    % TODO: CREATE AND ADD A DIAGRAM OF THE LEG KINEMATICS
    % 1. Add a schematic drawing of the leg showing
    %    L1, L2, L3 and Theta angles.
    % 2. Place it in 'src/paper/figures/kinematic_model.png'
    % 3. Uncomment the \includegraphics line below and
    %    delete the \rule placeholder!
    % \includegraphics[width=0.48\textwidth]{figures/kinematic_model.png}
    \rule{0.48\textwidth}{5cm} % Grey placeholder box
    \caption{The Denavit-Hartenberg (DH) geometric reference
    frame assigned to the three-degree-of-freedom quadruped limb.}
    \label{fig:kinematic_model}
\end{figure}

\subsection{Quintic Trajectory Formulation}

It is mathematically imperative that physical ground impact
velocities approach zero; violent mechanical shocks corrupt
internal IMU localization loops and induce structural fatigue.
Therefore, smooth, jerk-minimized foot trajectories are formulated
by employing overlapping piecewise quintic polynomials rather
than generic linear extrapolations.

For a specific foot transition evolving from resting coordinate
$A$ (lift-off) and traversing to coordinate $D$ (touch-down)
across constrained spanning time $T_{swing}$, an independent
trajectory locus $p(t)$ enforces continuous physical boundaries
in positional space, derivative velocity, and second-derivative
acceleration.

\begin{equation}
    p(t) = a_0 + a_1 t + a_2 t^2 + a_3 t^3 + a_4 t^4 + a_5 t^5
\end{equation}

By perfectly aligning adjacent spline limits exactly at
$t=0$ and $t=T_{swing}$, the iterative loops entirely preclude
unmodeled physical resonance. Overlapping standard trotting patterns
are generated efficiently via phase-shifting scalar arrays
(e.g., $0.50$ shifting applied to diagonal limb pairs).
